\documentclass[12pt]{article}
\usepackage{amsmath, amssymb}
\usepackage[margin=1in]{geometry}
\setlength{\parskip}{6pt}
\title{QUBO Formulation: Maximum Clique Problem}
\date{}
\begin{document}
\maketitle

\subsection*{Maximum Clique Problem}

\paragraph{Problem Statement.}
Let $G = (V, E)$ be an undirected graph with vertex set $V$ and edge set $E$. A clique is a subset $C \subseteq V$ such that every pair of distinct vertices in $C$ is connected by an edge in $E$. The Maximum Clique Problem seeks a clique $C^*$ of maximum cardinality, i.e., $|C^*| = \max \{ |C| : C \subseteq V \text{ and } \forall u,v \in C, \{u,v\} \in E \}$.

\paragraph{Decision Variables.}
For each vertex $i \in V$, we introduce a binary decision variable $x_i \in \{0,1\}$. The variable $x_i$ equals $1$ if vertex $i$ is included in the candidate clique, and $0$ otherwise.

\paragraph{QUBO Derivation.}
Step 1 --- The original objective is to maximize the number of selected vertices:
\begin{equation}
    \max \sum_{i \in V} x_i.
\end{equation}
Since QUBO solvers minimize energy, we negate the objective to obtain the minimization form:
\begin{equation}
    \min -\sum_{i \in V} x_i.
\end{equation}

Step 2 --- The feasibility condition requires that for any pair of vertices $\{i,j\}$ that are not connected by an edge (i.e., $\{i,j\} \notin E$), they cannot both be selected. This yields the inequality constraint:
\begin{equation}
    x_i + x_j \le 1, \quad \forall \{i,j\} \notin E.
\end{equation}

Step 3 --- To incorporate this inequality into the objective, we construct a penalty term. Because $x_i, x_j \in \{0,1\}$, the product $x_i x_j$ is $1$ if and only if both variables are $1$, and $0$ otherwise. Thus, the term $\lambda x_i x_j$ correctly penalizes only the infeasible configuration $(1,1)$ while leaving $(0,0)$, $(1,0)$, and $(0,1)$ unpenalized. We explicitly avoid the equality penalty $(x_i + x_j - 1)^2$, as it evaluates to $1$ for both $(0,0)$ and $(1,1)$, incorrectly penalizing valid configurations where neither vertex is selected. The total penalty contribution is:
\begin{equation}
    \lambda \sum_{\{i,j\} \notin E} x_i x_j.
\end{equation}

Step 4 --- Combining the objective and penalty terms gives the single QUBO Hamiltonian:
\begin{equation}
    H(\mathbf{x}) = -\sum_{i \in V} x_i + \lambda \sum_{\{i,j\} \notin E} x_i x_j.
\end{equation}

Step 5 --- Reading off the coefficients for the standard QUBO form $H(\mathbf{x}) = \sum_{i} Q_{i,i} x_i + \sum_{i<j} Q_{i,j} x_i x_j + c$, we obtain the following closed-form expressions for the matrix entries and offset:
\begin{align}
    Q_{i,i} &= -1, \quad \forall i \in V, \\
    Q_{i,j} &= \lambda, \quad \forall \{i,j\} \notin E, \\
    Q_{i,j} &= 0, \quad \forall \{i,j\} \in E, \\
    c &= 0.
\end{align}

\paragraph{Penalty Weight Justification.}
The penalty weight $\lambda$ must be chosen large enough to ensure that any violation of the clique constraints results in a strictly higher energy than the best possible feasible solution. Violating a single constraint $x_i + x_j \le 1$ corresponds to adding at most one extra vertex to the clique, which increases the objective sum by at most $1$. Therefore, any $\lambda > 1$ guarantees that the energy cost of the penalty outweighs the maximum possible gain from the objective. We set $\lambda = 2$ to satisfy this condition with a safe margin.

\paragraph{Correctness Argument.}
Minimizing $H(\mathbf{x})$ over $\{0,1\}^N$ recovers the maximum clique. For any feasible solution (a valid clique), all non-adjacent pairs have at least one variable equal to $0$, making every penalty term $x_i x_j$ equal to $0$. Thus, $H(\mathbf{x}) = -|C|$, and minimizing $H$ is equivalent to maximizing $|C|$. For any infeasible solution, at least one pair $\{i,j\} \notin E$ has $x_i = x_j = 1$, contributing at least $\lambda$ to the energy. Since $\lambda > 1$ and the maximum possible objective gain from adding a vertex is $1$, the energy of any infeasible configuration strictly exceeds that of the optimal feasible configuration. Consequently, the global minimum of $H$ must be feasible and correspond to a maximum clique.

\paragraph{Illustrative Example.}
Consider a graph with $N=3$ vertices $V=\{0,1,2\}$ and no edges, i.e., $E=\emptyset$. Here, every pair is a non-edge. Setting $\lambda=2$, the general formulas instantiate as follows. The objective terms are $-x_0 - x_1 - x_2$. The penalty terms are $2x_0x_1 + 2x_0x_2 + 2x_1x_2$. The resulting Hamiltonian is:
\begin{equation}
    H(\mathbf{x}) = -x_0 - x_1 - x_2 + 2x_0x_1 + 2x_0x_2 + 2x_1x_2.
\end{equation}
Evaluating $H$ over all $2^3=8$ binary assignments yields the minimum energy $-1$, achieved at configurations such as $(x_0,x_1,x_2) = (1,0,0)$, $(0,1,0)$, or $(0,0,1)$. These correspond to selecting exactly one vertex, which is the maximum clique size for an empty graph.

\end{document}